\Askhsh{1741}{4}
\begin{ylh}{1}
    \item 5.2. Παραλληλόγραμμα
    \item 5.6. Εφαρμογές στα τρίγωνα.
\end{ylh}
Δίνεται τρίγωνο $\mathrm{A}\mathrm{B}\Gamma$ και έστω $\mathrm{K}, \Lambda$ τα μέσα των πλευρών του $\mathrm{A}\mathrm{B}$ και $\mathrm{A}\Gamma$ αντίστοιχα.
\begin{alist}
    \item Θεωρούμε τυχαίο σημείο $\mathrm{M}$ στο εσωτερικό του τριγώνου και $\Delta, \mathrm{E}$ τα συμμετρικά του $\mathrm{M}$ ως προς $\mathrm{K}$ και $\Lambda$ αντίστοιχα. Να αποδείξετε ότι $\Delta\mathrm{E} \parallel \mathrm{B}\Gamma$. \monades{15}
    \item Στην περίπτωση που το σημείο $\mathrm{M}$ είναι το μέσο της πλευράς $\mathrm{B}\Gamma$, και $\Delta, \mathrm{E}$ τα συμμετρικά του $\mathrm{M}$ ως προς $\mathrm{K}$ και $\Lambda$ αντίστοιχα. Να αποδείξετε ότι τα σημεία $\Delta, \mathrm{A}$ και $\mathrm{E}$ είναι συνευθειακά. \monades{10}
\end{alist}
